\documentclass{report}
\usepackage{amsfonts, amsmath}
\newcommand\R{{\mathbb R}}
\newcommand\Q{{\mathbb Q}}
\newcommand\N{{\mathbb N}}
\newcommand\Or{{\mathcal O}}
\newcommand\ve{\varepsilon}
%\pagestyle{empty}
\begin{document}
\large
\centerline{\Huge Fisica Matematica II}
\renewcommand{\thefootnote}{ } 
\renewcommand{\thefootnote}{\arabic{footnote}} 
\vskip.1cm
\centerline{\Large Soluzioni esonero 09-06-2003}
\vskip1cm
\begin{enumerate}
\item Si cerca una soluzione del tipo $u(x,t)=xe^{-t}+v(x,t)$. Allora
\[
\begin{split}
&v_{xx}-v=v_t  \\
&v(x,0)=1-x;\; v(1,t)=0;\;v_x(0,t)=0.
\end{split}
\]
Si procede poi per separazione di variabili. Cio\`e
\[
v(x,t)=\sum_{n=0}^\infty c_n(t)\cos(n+\frac 12)\pi x
\]
che soddisfa le condizioni al contoro. Sostituendo si ha
\[
c_n'=-[1-(n+\frac 12)^2\pi^2]c_n;\quad c_n(0)=b_n
\]
dove
\[
b_n=2\int_0^1(1-x)\cos(n+\frac 12)\pi x=\frac{2}{(n+\frac 12)^2\pi^2}.
\]
Dunque
\[
c_n(t) =e^{-[1-(n+\frac 12)^2\pi^2]t}b_n.
\]
\item Si risolve per separazione di variabili nella forma
\[
u(x,t)=\sum_{n=1}^\infty c_n(t)\sin n\pi x.
\]
Sostituendo si ottiene il sistema
\[
\begin{split}
&c_n''=-n^2\pi^2c_n+a_n\\
&c_n(0)=c_n'(0)=0.
\end{split}
\]
Dove 
\[
a_n=\frac 1\pi\int_0^1 x\sin n\pi x\;dx.
\]
Le cui soluzioni sono
\[
c_n(t)=\frac{a_n}{n^2\pi^2}(1-\cos n\pi t).
\]
\item Si ha
\[
\hat f(\xi)=\frac{\alpha}{\alpha^2+\xi^2}.
\]
Da cui
\[
\lim_{\alpha\to 0}\int_\R \frac{\alpha}{\alpha^2+\xi^2} g(\xi)d\xi =
\lim_{\alpha\to 0}\int_\R\frac 1{1+z^2}g(\alpha z) dz=g(0).
\]
\item Poich\`e la velocit\`a di propagazione \`e limitata da 1, segue
che $u$ \`e a supporto compatto per ogni $t$. Dunque \`e facile
verificare che l'energia totale
\[
E(t)=\int_\R u_t^2+u_x^2
\]
\`e ben definita e costante (basta derivare rispetto al tempo e usare
il fatto che $u$ risolve l'equazione). Inoltre $e(0)=E(0)$, dunque
$e(t)\leq E(t)=E(0)=e(0)$. 
\end{enumerate}
\end{document}
